\subsection{Isotropy group of a coequalizer}

Let G and H be groupoids; let \(\varphi\) and \(\psi\) morphisms of groupoids from H to G. The purpose of this section is to summarize the computation of the isotropy groups of the cohomotoper (II, p. 193, prop. 6) and the comparison of the isotropy groups of the cohomotoper and of the coequalizer made in the preceding section, in order to deduce from this the computation of the isotropy groups of the coequalizer Coeg( \(\varphi\) , \(\psi\) ).

Denote by \(\alpha\colon\mathrm{G}\to\mathrm{Coh}(\varphi,\psi)\) , \(\beta\colon\mathrm{Coh}(\varphi,\psi)\to\mathrm{Coeg}(\varphi,\psi)\) and \(\gamma=\beta\circ\alpha\colon\mathrm{G}\to\mathrm{Coeg}(\varphi,\psi)\) the canonical groupoid morphisms. Denote \(h\colon\mathrm{Som}(\mathrm{H})\to\mathrm{Fl}(\mathrm{Coh}(\varphi,\psi))\) the canonical homotopy; recall that the set of vertices of \(\mathrm{Coh}(\varphi,\psi)\) is equal to \(\mathrm{Som}(\mathrm{G})\) .

Denote by \(\Gamma\) the quiver (Som(G), Som(H), Som(φ), Som(ψ)), denote then \(\varphi_0\) and \(\psi_0\) the maps induced by \(\varphi\) and \(\psi\) by passing to orbits and \(\Gamma_0\) the frame (Orb(G), Orb(H), \(\varphi_0, \psi_0\) ) of the pair ( \(\varphi, \psi\) ). We shall assume that the quiver \(\Gamma_0\) is connected and its vertex set is nonempty, which amounts to assuming that the groupoids Coh(φ, ψ) and Coeg(φ, ψ) are transitive.

Recall (cf. II, p. 192, definition 4) that a base equipment consists of the data of a family \((a, b, c_1, c_2, \mathrm{T}, i_0)\) where: for all \(i \in \mathrm{Orb}(\mathrm{G})\) , \(a(i)\) is a vertex in the orbit \(i\) of \(\mathrm{G}\) ; for every \(j \in \mathrm{Orb}(\mathrm{H})\) , \(b(j)\) is a vertex in the orbit \(j\) of \(\mathrm{H}\) , \(c_1(j)\) and \(c_2(j)\) are arrows from \(\mathrm{G}\) respectively \(\varphi(b(j))\) to \(a(\varphi_0(j))\) and \(\psi(b(j))\) to \(a(\psi_0(j))\) ; \(\mathrm{T}\) is a subquiver of \(\Gamma_0\) whose associated tree is a maximal tree of the graph \(\widetilde{\Gamma}_0\) ; finally, \(i_0\) is an orbit of \(\mathrm{G}\) and we set \(a_0 = a(i_0)\) .

We then define a quiver morphism \(\tau_0\) from \(\Gamma_0\) to \(\mathrm{Coh}(\varphi, \psi)\) such that \(\tau_0(i) = \beta(a(i))\) and \(\tau_0(j) = \alpha(c_1(j))^{-1}h(b(j))\alpha(c_2(j))\) for \(i \in \mathrm{Orb}(\mathrm{G})\) and \(j \in \mathrm{Orb}(\mathrm{H})\) . If \(i\) belongs to \(\mathrm{Orb}(\mathrm{G})\) , let \(d_i\) be the unique class of paths in the graph \(\widetilde{\mathrm{T}}\) connecting \(i_0\) to \(i\) and denote by \(\delta_i\) its image in \(\mathrm{Coh}(\varphi, \psi)\) by the groupoid morphism \(\tau_0\) .

From these data, prop. 6 (II, p. 193) yields a surjective homomorphism.

\[
\Lambda \colon \bigl (\underset {i \in \pi_ {0} (\mathrm{G})} {*} \mathrm{G} _ {a (i)} \bigr) * \mathrm{F} (\operatorname{Orb} (\mathrm{H})) \to \operatorname{Coh} (\varphi , \psi) _ {a (i _ {0})}
\]

and describes generators of its kernel, thereby providing a presentation of the group \(\mathrm{Coh}(\varphi,\psi)_{a_{0}}\) .

The set of loops of length \(\geqslant 1\) in the graph \(\widetilde{\Gamma}\) associated with \(\Gamma\) is identified with the set \(\Omega(\widetilde{\Gamma})\) of sequences \((z_1, \ldots, z_n)\) of arrows of \(\widetilde{\Gamma}\) , indexed by \(\mathbf{Z}/n\mathbf{Z}\) , where \(n\) runs through the set of integers \(\geqslant 1\) , such that for every \(k \in \mathbf{Z}/n\mathbf{Z}\) , the term of \(z_k\) is the origin of \(z_{k+1}\) . Let \(\mathbf{Z}\) be a subset of \(\Omega(\widetilde{\Gamma})\) such that the conjugacy classes \(c(z)\) , for \(z \in \mathbf{Z}\) , generate the kernel of the surjective homomorphism \(\beta_{a_0}\) . The homomorphism \(\beta_{a_0} \circ \Lambda\) is surjective; to deduce its kernel from this, that is, a presentation of the group \(\operatorname{Coeg}(\varphi, \psi)_{\beta(a_0)}\) , it remains to choose, for every \(z \in \mathbf{Z}\) , an element \(\mathrm{C}(z)\) in the group \(\left( \begin{array}{cc} * & \mathrm{G}_{a(i)} \\ i \in \mathrm{Orb}(\mathrm{G}) & \end{array} \right) * \mathrm{F}(\pi_0(\mathrm{H}))\) such that \(\Lambda(\mathrm{C}(z))\) belongs to the conjugacy class \(c(z)\) .

So let \(z = (z_{1}, \ldots, z_{n})\) be an element of \(\Omega(\tilde{\Gamma})\) ; set \(z_{k} = (y_{k}, \varepsilon_{k})\) , where \(y_{k} \in \mathrm{Fl}(\Gamma) = \mathrm{Som}(\mathrm{H})\) and \(\varepsilon_{k} \in \{\pm 1\}\) . By definition, \(c(z)\) is the conjugacy class of the element

\[
g h (y _ {1}) ^ {\varepsilon_ {1}} \dots h (y _ {n}) ^ {\varepsilon_ {n}} g ^ {- 1}
\]

of the group \(\mathrm{Coh}(\varphi,\psi)_{a_{0}}\) , where g is an arbitrary arrow in \(\mathrm{Coh}(\varphi,\psi)\) connecting \(a_{0}\) to the origin of \(z_{1}\) . For \(k \in Z/nZ\) , let \(j_{k}\) be the orbit of \(y_{k}\) in H and choose an arrow \(f_{k}\) of H connecting \(y_{k}\) to the vertex \(b(j_{k})\). By definition of a homotopy, we then have the relation

\[
h (y _ {k}) (\alpha \circ \psi) (f _ {k}) = (\alpha \circ \varphi) (f _ {k}) h (b (j _ {k}))
\]

in the groupoid \(\mathrm{Coh}(\varphi,\psi)\) . Consequently, using the definition of the quiver morphism \(\tau_{0}\) , we have

\[
\begin{array}{l} h (y _ {k}) = (\alpha \circ \varphi) (f _ {k}) \cdot h (b (j _ {k})) \cdot (\alpha \circ \psi) (f _ {k}) ^ {- 1} \\ \quad = (\alpha \circ \varphi) (f _ {k}) \cdot (\alpha \circ c _ {1}) (j _ {k}) \cdot \tau_ {0} (j _ {k}) \cdot (\alpha \circ c _ {2}) (j _ {k}) ^ {- 1} \cdot (\alpha \circ \psi) (f _ {k}) ^ {- 1} \\ \quad = (\alpha \circ \varphi) (f _ {k}) \cdot (\alpha \circ c _ {1}) (j _ {k}) \cdot \delta_ {\varphi_ {0} (j _ {k})} ^ {- 1} \cdot \\ \quad \cdot \delta_ {\varphi_ {0} (j _ {k})} \cdot \tau_ {0} (j _ {k}) \cdot \delta_ {\psi_ {0} (j _ {k})} ^ {- 1} \cdot \\ \quad \cdot \delta_ {\psi_ {0} (j _ {k})} \cdot (\alpha \circ c _ {2}) (j _ {k}) ^ {- 1} \cdot (\alpha \circ \psi) (f _ {k}) ^ {- 1} \\ = u _ {k} \Lambda (j _ {k}) v _ {k}, \end{array}
\]

where we have set

\[
u _ {k} = \alpha (\varphi (f _ {k}) c _ {1} (j _ {k})) \cdot \delta_ {\varphi_ {0} (j _ {k})} ^ {- 1} \quad \text { and } \quad v _ {k} = \delta_ {\psi_ {0} (j _ {k})} \cdot \alpha (\psi (f _ {k}) c _ {2} (j _ {k})) ^ {- 1}.
\]

For every element \(k \in \mathbf{Z}/n\mathbf{Z}\) , define arrows \(\widetilde{u}_{k}, \widetilde{v}_{k}\) in G by

\[
h (y _ {k}) ^ {\varepsilon_ {k}} = \widetilde {u} _ {k} \Lambda (j _ {k}) ^ {\varepsilon_ {k}} \widetilde {v} _ {k},
\]

so that

\[
(\widetilde {u} _ {k}, \widetilde {v} _ {k}) = \left\{ \begin{array}{l l} (u _ {k}, v _ {k}) & \text { if } \varepsilon_ {k} = 1; \\ (v _ {k} ^ {- 1}, u _ {k} ^ {- 1}) & \text { if } \varepsilon_ {k} = - 1. \end{array} \right.
\]

Denote by \(x_{k}\) the origin of the arrow \(h(y_{k})^{\varepsilon_{k}}\) ; its target is then \(x_{k+1}\) ; let \(i_{k}\) be the orbit of \(x_{k}\) . Let us then define a loop \(\lambda_{k}(z)\) at \(a(i_{k})\) in the groupoid G by the formula

\[
\lambda_ {k} (z) = \left\{ \begin{array}{l l} c _ {2} (j _ {k - 1}) ^ {- 1} \psi (f _ {k - 1}) ^ {- 1} \varphi (f _ {k}) c _ {1} (j _ {k}) & \text {if} (\varepsilon_ {k - 1}, \varepsilon_ {k}) = (1, 1); \\ c _ {2} (j _ {k - 1}) ^ {- 1} \psi (f _ {k - 1}) ^ {- 1} \psi (f _ {k}) c _ {2} (j _ {k}) & \text {if} (\varepsilon_ {k - 1}, \varepsilon_ {k}) = (1, - 1); \\ c _ {1} (j _ {k - 1}) ^ {- 1} \varphi (f _ {k - 1}) ^ {- 1} \varphi (f _ {k}) c _ {1} (j _ {k}) & \text {if} (\varepsilon_ {k - 1}, \varepsilon_ {k}) = (- 1, 1); \\ c _ {1} (j _ {k - 1}) ^ {- 1} \varphi (f _ {k - 1}) ^ {- 1} \psi (f _ {k}) c _ {2} (j _ {k}) & \text {if} (\varepsilon_ {k - 1}, \varepsilon_ {k}) = (- 1, - 1) \end{array} \right.\tag{1}
\]

By construction, we have

\[
\Lambda (\lambda_ {k} (z)) = \widetilde {v} _ {k - 1} \widetilde {u} _ {k},
\]

so that the image under the homomorphism \(\Lambda\) of the element

\[
\mathrm{C} (z) = \lambda_ {1} (z) (j _ {1}) ^ {\varepsilon_ {1}} \lambda_ {2} (z) (j _ {2}) ^ {\varepsilon_ {2}} \dots \lambda_ {n} (z) (j _ {n}) ^ {\varepsilon_ {n}}\tag{2}
\]

belongs to the conjugacy class \(c(z)\) .

DEFINITION 3. --- Let G and H be groupoids; let \(\varphi\) and \(\psi\) morphisms of groupoids from H to G. We assume that the groupoid Coeg \((\varphi, \psi)\) is transitive. Let \((a, b, c_1, c_2, T, i_0)\) be a basic equipment of the pair \((\varphi, \psi)\) .

One calls complementary equipment the data of a subset Z of \(\Omega(\widetilde{\Gamma})\) such that the conjugacy classes \(c(\mathbf{z})\) for \(\mathbf{z} \in \mathbf{Z}\) generate the kernel of the homomorphism \(\beta_{a(i_0)}\) and, for every element \(\mathbf{z} = ((y_1, \varepsilon_1), \ldots, (y_n, \varepsilon_n))\) of Z, a sequence \(f(\mathbf{z}) = (f_1, \ldots, f_n)\) of arrows in H such that \(f_k\) connects \(y_k\) to the vertex \(b(j_k)\) , \(j_k\) denoting the orbit of \(y_k\) in H.

A complete equipment of the pair \((\varphi,\psi)\) is the data of a base equipment and a complementary equipment.

PROPOSITION 5. --- Let G and H be groupoids; let \(\varphi\) and \(\psi\) be morphisms of groupoids from H to G. Suppose that the groupoid \(\operatorname{Coeg}(\varphi, \psi)\) is transitive. Denote by \(\gamma\colon \mathrm{G} \to \mathrm{Coeg}(\varphi, \psi)\) the canonical groupoid morphism.

We equip the pair \((\varphi, \psi)\) with a complete equipment

\[
(a, b, c _ {1}, c _ {2}, \mathrm{T}, i _ {0}, \mathrm{Z}, (f (\mathbf {z})) _ {\mathbf {z} \in \mathrm{Z}}).
\]

For \(j \in \text{Orb(H)}\) , let \(\varphi_j \colon \text{H}_{b(j)} \to \text{G}_{a(\varphi_0(j))}\) and \(\psi_j \colon \text{H}_{b(j)} \to \text{G}_{a(\psi_0(j))}\) be the group homomorphisms \(\text{Int}(c_1(j))^{-1} \circ \varphi_{b(j)}\) and \(\text{Int}(c_2(j))^{-1} \circ \psi_{b(j)}\) respectively, where \(\varphi_0\) and \(\psi_0\) denote the canonical maps induced by \(\varphi\) and \(\psi\) by passing to orbits. Let \(\tau\) be the morphism of the frame \(\Gamma_0\) of the pair \((\varphi, \psi)\) into \(\text{Coeg}(\varphi, \psi)\) defined by \(\text{Som}(\tau)(i) = \gamma(a(i))\) if \(i \in \text{Orb(G)}\) and such that \(\text{Fl}(\tau)(j)\) is the path \(\gamma(c_1(j))^{-1} \gamma(c_2(j))\) in \(\text{Coeg}(\varphi, \psi)\). If \(i\) is an orbit of G, let \(c_i\) be the unique class of paths in \(\widetilde{\text{T}}\) connecting \(i_0\) to \(i\), and set \(\delta_i = \widetilde{\tau}(c_i)\) , where \(\widetilde{\tau} \colon \text{Grp(G)} \to \text{Coeg}(\varphi, \psi)\) is the canonical groupoid morphism induced by \(\tau\) .

There then exists a unique group homomorphism

\[
\lambda \colon \bigl ( \begin{array}{c} * \\ i \in \operatorname{Orb} (\mathrm{G}) \end{array} \mathrm{G} _ {a (i)} \bigr) * \mathrm{F} (\operatorname{Orb} (\mathrm{H})) \to \operatorname{Coeg} (\varphi , \psi) _ {\gamma (a (i _ {0}))}
\]

such that

\[
\begin{array}{l l} \lambda (f) = \delta_ {i} \gamma_ {a (i)} (f) \delta_ {i} ^ {- 1} & \text { for } i \in \mathrm{Orb(G)} \text { and } f \in \mathrm{G} _ {a (i)}, \\ \lambda (j) = \delta_ {\varphi_ {0} (j)} \tau (j) \delta_ {\psi_ {0} (j)} ^ {- 1} & \text { for } j \in \mathrm{Orb(H)}. \end{array}
\]

The homomorphism \(\lambda\) is surjective; its kernel is the smallest normal subgroup containing the following elements:

\[
\left(\mathrm{R} _ {1}\right) \quad r _ {1} (j) = j \quad \text { for } j \text { in } \mathrm{Fl} (\mathrm{T});
\]

\[
(\mathrm{R} _ {2}) \quad r _ {2} (j, f) = \varphi_ {j} (f) j \psi_ {j} (f) ^ {- 1} j ^ {- 1}
\]

\[
\text { for } j \in \operatorname{Orb} (\mathrm{H}) \text { and } f \in \mathrm{H} _ {b (j)};\tag{\((R_{3})\)}
\]

\[
r _ {3} (z) = \lambda_ {1} (z) j _ {1} ^ {\varepsilon_ {1}} \lambda_ {2} (z) j _ {2} ^ {\varepsilon_ {2}} \dots \lambda_ {n} (z) j _ {n} ^ {\varepsilon_ {n}}
\]

\[
\text { for } z = ((y _ {1}, \varepsilon_ {1}), \dots , (y _ {n}, \varepsilon_ {n})) \in \mathbf {Z},
\]

where the loops \(\lambda_{i}(z)\) are defined by formula (1), p. 208.

The existence and uniqueness of such a morphism follows from the universal property of the free product of a family of groups (A, I, p. 85, prop. 8). Denote by \(\alpha\colon G\to\mathrm{Coh}(\varphi,\psi)\) and \(\beta\colon\mathrm{Coh}(\varphi,\psi)\to\mathrm{Coeg}(\varphi,\psi)\) the canonical morphisms, so that \(\gamma=\beta\circ\alpha\). Let also \(h\colon\mathrm{Som}(H)\to\mathrm{Fl}(\mathrm{Coh}(\varphi,\psi))\) be the canonical homotopy relating \(\alpha\circ\varphi\) to \(\alpha\circ\psi\) . Since \(\beta\) is the groupoid morphism obtained by contracting the arrows of the image of \(h\) , we have

\[
\operatorname{Fl} (\tau) (j) = \gamma (c _ {1} (j)) ^ {- 1} \gamma (c _ {2} (j)) = \gamma (c _ {1} (j) ^ {- 1}) \beta (h (j)) \gamma (c _ {2} (j)) = \beta (\tau_ {0} (j))
\]

in \(\operatorname{Coeg}(\varphi,\psi)\) , where \(\tau_{0}\colon\Gamma_{0}\to\mathrm{Coh}(\varphi,\psi)\) denotes the quiver morphism induced by the base equipment \((a,b,c_{1},c_{2},\mathrm{T},i_{0})\) . Consequently, the homomorphism \(\lambda\) is the composite of the homomorphism \(\Lambda\) defined in Prop. 6 (II, p. 193) and of the surjective homomorphism \(\beta_{a(i_0)}\) . It is in particular surjective.

Let \(z \in \mathbf{Z}\). By construction, \(\Lambda(r_3(z))\) belongs to the conjugacy class \(c(z)\) . By definition of a complementary equipment, the kernel of the homomorphism \(\beta_{a(i_0)}\) is therefore the smallest normal subgroup of \(\mathrm{Coh}(\varphi, \psi)_{a(i_0)}\) containing the elements \(\Lambda(r_3(z))\) for \(z \in \mathbf{Z}\) . Since the homomorphism \(\Lambda\) is surjective, the kernel of the homomorphism \(\lambda = \beta_{a(i_0)} \circ \Lambda\) is therefore the smallest normal subgroup of the group \(\left( *_{i \in \mathrm{Orb}(G)} G_{a(i)} \right) * F(\pi_0(H))\) containing the generators of the kernel of \(\Lambda\) given by the formulas \((R_1), (R_2)\) , as well as the elements defined by the formulas \((R_3)\) . The proposition is thus proved.

